\begin{proof}

\begin{enumerate}[label={$\langle1\rangle$\arabic*.},leftmargin=*]
\item \ref{push:src:E.2.12}, \ref{push:src:E.2.14} and \ref{push:src:E.2.15} exclude every degenerate
limit of the normalizing projective transformations.
They therefore extend as invertible projective
transformations, and the three labelled critical
values remain distinct.  \ref{push:src:E.2.13} proves the stated
finite-fibre multiplicities.

\item Normalize the three sections to $0,1,\infty$.
Over $B^\circ=\mathbb P^1\setminus\{0,1,\infty\}$,
the relative map is a proper submersion over
$B^\circ\times\Delta$.  Hence its integral
first-homology monodromies on the central fibre
are the same matrices $A_1,A_2,A_0$ as in
\ref{push:src:E.2.7}.  The common invariant space of the dual
matrices $A_1^t,A_2^t$ has dimension one:
in the displayed marking it is generated by
the covector $\gamma$.  This is obtained by
solving their simultaneous fixed equations;
the exact calculation is recorded in
\path{code/e_push_interior_cusp.py}.

\item Let $r$ be the total number of irreducible
components of the three reduced critical
fibres, and let $X^\circ$ be their complement
in $X_0$.  Alexander duality on $S^6$ gives
\[
 b_1(X^\circ;\mathbb Q)=r.
\]
The degree-one Leray sequence over the
punctured projective line, whose second
cohomology vanishes, gives instead
\[
 b_1(X^\circ;\mathbb Q)
 =b_1(B^\circ;\mathbb Q)
   +\dim H^1(T;\mathbb Q)^{\rm mon}
 =2+1=3.
\]
There are three nonempty critical fibres,
so each has exactly one reduced component.

\item Write the cusp fibre as $dS$ with $d\ge1$.
We prove $d=1$ by a direct meridian argument.
The smooth fibration over $B^\circ$ has
connected fibres, hence
$\pi_1(X^\circ)\to\pi_1(B^\circ)$ is
surjective.  The inclusion
$X^\circ\hookrightarrow X_0$ is surjective
on fundamental groups, and its kernel is
normally generated by meridians about the
three reduced divisor components.  To see
the latter assertion, make a null-homotopy
disc transverse to their smooth loci and
disjoint from their singular loci, which
have real codimension at least four.
Removing small discs about the finitely
many intersections expresses its boundary
as a product of conjugate meridians.

\item A meridian of a multiplicity-$a$ component
maps to the $a$th power of the corresponding
base meridian, since locally the base
function is a unit times the $a$th power
of a reduced local equation.  Consequently
there is a surjection
\[
 \pi_1(X_0)\longrightarrow
 \left\langle a,b,c\ \middle|\
 a^m=b^n=c^d=abc=1\right\rangle.
 \tag{E.2.16.1}
\]
This is the usual orbifold fundamental
group, but the surjection has just been
proved without assuming a stack
factorization.  Its source is trivial.

\item For clarity the group on the right is
nontrivial for every $d\ge2$ here.
Except for $(m,n,d)=(3,4,2)$, one has
\[
 \frac1m+\frac1n+\frac1d<1.
\]
Indeed for $d\ge3$ the sum is at most
$11/12$; for $d=2$, increasing either
$m=3k$ or $n=4l$ above $3,4$ makes it
at most $23/24$.  Put
$\alpha=\pi/m$, $\beta=\pi/n$,
$\theta=\pi/d$, and choose $\ell>0$ with
\[
 \cosh\ell=
 \frac{\cos\theta+\cos\alpha\cos\beta}
      {\sin\alpha\sin\beta}>1.
\]
Let
\[
 R_\alpha=\begin{pmatrix}
 \cos\alpha&\sin\alpha\\
 -\sin\alpha&\cos\alpha
 \end{pmatrix},\qquad
 D_\ell=\operatorname{diag}(e^{\ell/2},e^{-\ell/2}).
\]
In $\operatorname{PSL}_2(\mathbb R)$ the
elements represented by $R_\alpha$ and
$D_\ell R_\beta D_\ell^{-1}$ have orders
$m,n$, and their product has trace
$-2\cos\theta$ and projective order $d$.
Taking its inverse for $c$ gives a
nontrivial representation of (E.2.16.1).
For the exceptional triple, the
permutations $a=(123)$ and $b=(1324)$
have orders three and four, and
$ab=(24)$ has order two, giving a
nontrivial representation in $S_4$.
Thus $d\ge2$ contradicts (E.2.16.1),
and the cusp is reduced.

\item The ordinary torus fibres have Euler
characteristic zero.  The finite reduced
fibres are smooth finite free quotients
of two-tori by \ref{push:src:E.2.13}, so also have
Euler characteristic zero.  Additivity
over the three singular values therefore
gives $\chi_{\rm top}(S)=\chi(X_0)=2$.
Finally, the valid irreducible-fibre
degree-one invariant-cycle statement
\cite[Corollary 3.7, arXiv version~1]{HondaViaclovsky2025D}
gives
\[
 b_1(S;\mathbb Q)=
 \dim H^1(T;\mathbb Q)^{A_0}=2.
\]
Its irreducibility hypothesis has just
been proved.  This use is separate from
the false unrestricted equality
converse audited in \SeriesRef{push:src:B.4.17}.

\end{enumerate}
\end{proof}
